\documentclass[11pt, letterpaper]{article}

% --- Idioma y codificación ---
\usepackage[spanish]{babel}
\usepackage[utf8]{inputenc}
\usepackage[T1]{fontenc}
\usepackage{lmodern}

% --- Matemáticas ---
\usepackage{amsmath, amssymb, amsthm}

% --- Formato ---
\usepackage[margin=2.5cm]{geometry}
\usepackage{parskip}
\usepackage{booktabs}
\usepackage{listings}
\usepackage{hyperref}
\hypersetup{colorlinks=true, linkcolor=blue!70!black, urlcolor=blue!70!black}

% --- Código ---
\lstset{
  language=C, basicstyle=\ttfamily\small,
  keywordstyle=\bfseries, commentstyle=\itshape,
  numbers=left, numberstyle=\tiny, frame=single,
  breaklines=true
}

% --- Entornos de teoremas ---
\theoremstyle{definition}
\newtheorem{definition}{Definición}[section]
\theoremstyle{plain}
\newtheorem{theorem}{Teorema}[section]
\newtheorem{lemma}[theorem]{Lema}
\newtheorem{corollary}[theorem]{Corolario}
\theoremstyle{remark}
\newtheorem{remark}{Observación}[section]

\title{%
  T02 --- Demostración formal:\\[4pt]
  La inserción preserva la propiedad BST
}
\author{Curso Linux--C--Rust --- B15 Estructuras de datos en C}
\date{}

\begin{document}
\maketitle
\tableofcontents

% ===================================================================
\section{Definiciones formales}
% ===================================================================

\begin{definition}[Propiedad BST]
Un árbol binario $T$ satisface la propiedad BST si para todo nodo $x$
en $T$:
\[
  \forall\, y \in \text{sub\_izq}(x):\; \text{key}(y) < \text{key}(x)
  \qquad\text{y}\qquad
  \forall\, z \in \text{sub\_der}(x):\; \text{key}(z) > \text{key}(x)
\]
\end{definition}

\begin{definition}[Rango válido]
Cada nodo $x$ tiene un rango $(lo, hi)$ tal que $lo < \text{key}(x) < hi$.
Los subárboles heredan rangos restringidos:
\begin{itemize}
  \item Hijo izquierdo: rango $(lo, \text{key}(x))$.
  \item Hijo derecho: rango $(\text{key}(x), hi)$.
\end{itemize}
La raíz tiene rango $(-\infty, +\infty)$.
\end{definition}

\begin{definition}[BST válido]
$T$ es un BST válido si $T$ es vacío (NULL), o $T$ satisface la
propiedad BST y el recorrido inorden produce una secuencia
estrictamente creciente.
\end{definition}

Asumimos claves únicas: la inserción ignora valores ya presentes.

% ===================================================================
\section{Demostración 1: inserción iterativa (invariante de bucle)}
% ===================================================================

Consideramos la inserción iterativa con puntero a puntero:

\begin{lstlisting}
void bst_insert_pp(TreeNode **root, int value) {
    while (*root) {
        if (value < (*root)->data)
            root = &(*root)->left;
        else if (value > (*root)->data)
            root = &(*root)->right;
        else
            return;  // duplicado
    }
    *root = create_node(value);
}
\end{lstlisting}

\begin{theorem}\label{thm:iterative}
Si el árbol $T$ es un BST válido antes de llamar a
\texttt{bst\_insert\_pp(\&T, v)}, entonces es un BST válido después.
\end{theorem}

\begin{definition}[Invariante de bucle]
\begin{align}
  \textbf{(I1)}&\quad \text{\texttt{*root} es un subárbol del BST original } T. \notag\\
  \textbf{(I2)}&\quad lo < v < hi, \text{ donde } (lo, hi) \text{ es el rango válido de la posición de \texttt{root}.} \notag\\
  \textbf{(I3)}&\quad \text{Los nodos del BST original no han sido modificados.} \notag
\end{align}
\end{definition}

\begin{proof}[Prueba del Teorema~\ref{thm:iterative}]
Probamos el invariante en tres etapas.

\textbf{Inicialización.} Al entrar, \texttt{root = \&T}.
\begin{itemize}
  \item (I1): \texttt{*root} $= T$, el árbol completo. \checkmark
  \item (I2): Rango de la raíz $= (-\infty, +\infty)$, y $-\infty < v < +\infty$. \checkmark
  \item (I3): Ninguna operación ejecutada. \checkmark
\end{itemize}

\textbf{Mantenimiento.} Supongamos $\text{INV}(\texttt{root}, lo, hi)$
con $\texttt{*root} \neq \text{NULL}$. Sea $x = \texttt{*root}$ con
clave $k$.

\emph{Caso A: $v < k$.} Se ejecuta \texttt{root = \&(x->left)}.
Nuevo rango: $(lo, k)$.
\begin{itemize}
  \item (I1): \texttt{*root} es \texttt{x->left}, subárbol del BST. \checkmark
  \item (I2): $lo < v$ (hipótesis) y $v < k$, así que $lo < v < k$. \checkmark
  \item (I3): Solo se reasigna el puntero local. \checkmark
\end{itemize}

\emph{Caso B: $v > k$.} Simétrico. Nuevo rango: $(k, hi)$, con
$k < v < hi$. \checkmark

\emph{Caso C: $v = k$.} Se retorna sin modificar el árbol. \checkmark

\textbf{Terminación.} El bucle termina cuando $\texttt{*root} = \text{NULL}$.
Por (I2), $lo < v < hi$. Se ejecuta $\texttt{*root} = \text{create\_node}(v)$,
creando una hoja con clave $v$.

El resultado es un BST válido porque:
\begin{enumerate}
  \item La hoja $v$ tiene subárboles vacíos: satisface la propiedad BST
    trivialmente.
  \item El padre de $v$: si $v$ es hijo izquierdo, $v < \text{key}(\text{padre})$;
    si derecho, $v > \text{key}(\text{padre})$. Correcto por (I2).
  \item Los ancestros: por (I2), $v$ está dentro del rango heredado de
    cada ancestro. Por (I3), sus relaciones de orden no cambian.
  \item Nodos fuera del camino: no se ven afectados.
\end{enumerate}
\end{proof}

\begin{lemma}[Terminación del bucle]\label{lem:loop-term}
El bucle termina en a lo sumo $h + 1$ iteraciones, donde $h$ es la
altura del árbol.
\end{lemma}

\begin{proof}
En cada iteración, \texttt{root} desciende un nivel (profundidad
$d \to d+1$). La profundidad máxima es $h$; tras $h$ descensos se
alcanza una hoja cuyo hijo es NULL.
\end{proof}

% ===================================================================
\section{Demostración 2: inserción recursiva (inducción estructural)}
% ===================================================================

\begin{lstlisting}
TreeNode *bst_insert(TreeNode *root, int value) {
    if (!root) return create_node(value);
    if (value < root->data)
        root->left = bst_insert(root->left, value);
    else if (value > root->data)
        root->right = bst_insert(root->right, value);
    return root;
}
\end{lstlisting}

\begin{theorem}\label{thm:recursive}
Si $T$ es un BST válido con rango $(lo, hi)$ y $lo < v < hi$, entonces
\texttt{bst\_insert(T, v)} retorna un BST válido con rango $(lo, hi)$.
\end{theorem}

\begin{proof}
Por inducción estructural sobre $T$.

\textbf{Base ($T = \text{NULL}$).} Se retorna una hoja con clave $v$.
Sin hijos, satisface la propiedad BST trivialmente. El rango
$(lo, hi)$ con $lo < v < hi$ se cumple. \checkmark

\textbf{Paso inductivo ($T$ no vacío).} Sea $r$ la raíz con clave $k$
y rango $(lo, hi)$.

\emph{Hipótesis inductiva}: para cualquier subárbol $S$ de $T$ con
rango $(lo', hi')$ y $lo' < v < hi'$, \texttt{bst\_insert(S, v)}
retorna un BST válido con rango $(lo', hi')$.

\emph{Caso $v < k$:} Se llama \texttt{bst\_insert(root->left, v)}.
El subárbol izquierdo tiene rango $(lo, k)$. Como $lo < v$ y $v < k$:
$lo < v < k$. Por hipótesis inductiva, el resultado es un BST válido
con rango $(lo, k)$. El subárbol derecho no cambia. La raíz $r$ sigue
satisfaciendo la propiedad BST. \checkmark

\emph{Caso $v > k$:} Simétrico, con rango $(k, hi)$. \checkmark

\emph{Caso $v = k$:} Se retorna \texttt{root} sin modificación. \checkmark
\end{proof}

% ===================================================================
\section{Demostración 3: inorden correcto tras la inserción}
% ===================================================================

\begin{theorem}\label{thm:inorder}
Sea $T$ un BST con recorrido inorden $[a_1, a_2, \ldots, a_n]$.
Después de insertar $v \notin T$, el inorden es
$[a_1, \ldots, a_j, v, a_{j+1}, \ldots, a_n]$ donde $a_j < v < a_{j+1}$.
\end{theorem}

\begin{proof}
La inserción coloca $v$ como hoja. En el recorrido inorden, una hoja
aparece entre su antecesor y sucesor inorden.

En el camino de la raíz a la hoja $v$:
\begin{itemize}
  \item Si se desciende a la izquierda en un nodo $x$ (porque $v < \text{key}(x)$),
    $x$ es candidato a \emph{sucesor} inorden de $v$.
  \item Si se desciende a la derecha ($v > \text{key}(x)$), $x$ es
    candidato a \emph{antecesor} inorden de $v$.
\end{itemize}

El antecesor inorden es el último nodo donde se fue a la derecha
(el más profundo con $k < v$); el sucesor es el último donde se fue a
la izquierda.

Como $v$ se inserta como hoja sin hijos, no altera el orden relativo
de los nodos existentes. El inorden del nuevo árbol es el anterior con
$v$ insertado en su posición ordenada.
\end{proof}

% ===================================================================
\section{Resumen de resultados}
% ===================================================================

\begin{table}[h]
\centering
\begin{tabular}{llp{6cm}}
\toprule
\textbf{Resultado} & \textbf{Técnica} & \textbf{Componentes} \\
\midrule
Inserción iterativa preserva BST
  & Invariante de bucle
  & Inicialización, mantenimiento (3 casos), terminación \\
Inserción recursiva preserva BST
  & Inducción estructural
  & Base (NULL), paso (3 casos) \\
Inorden correcto tras inserción
  & Antecesor/sucesor inorden
  & Propiedad del camino raíz $\to$ hoja \\
Terminación del bucle
  & Función de decrecimiento
  & Profundidad creciente, acotada por $h$ \\
\bottomrule
\end{tabular}
\end{table}

\end{document}
